\documentclass[11pt]{article}

\usepackage[a4paper,margin=1in]{geometry}
\usepackage{amsmath,amssymb,amsthm,mathtools}
\usepackage[T1]{fontenc}
\usepackage{lmodern}
\usepackage{microtype}
\usepackage{hyperref}

\newtheorem{theorem}{Theorem}
\newtheorem{lemma}[theorem]{Lemma}
\newtheorem{proposition}[theorem]{Proposition}
\newtheorem{corollary}[theorem]{Corollary}
\newtheorem{remark}[theorem]{Remark}
\newtheorem{definition}[theorem]{Definition}

\title{Connectedness thresholds for a tridiagonal Toeplitz pseudospectrum}
\author{}
\date{}

\begin{document}
\maketitle

Put
\[
q=\sqrt a,\qquad
D_n=\operatorname{diag}(1,q,\ldots,q^{n-1}),
\qquad
T_n=\operatorname{tridiag}(1,0,1).
\]
Then
\[
D_n^{-1}A_n(a)D_n=qT_n.
\]
Consequently,
\[
\sigma(A_n(a))
=
\left\{
\lambda_{n,j}=2q\cos\frac{j\pi}{n+1}:1\le j\le n
\right\},
\]
where
\[
\lambda_{n,1}>\lambda_{n,2}>\cdots>\lambda_{n,n}.
\]
Set
\[
\rho_n=2q\cos\frac{\pi}{n+1}
\]
and
\[
f_n(z)=s_{\min}(zI_n-A_n(a)).
\]

We first record the finite drifted-Jacobi oscillation lemma that will be used below.

\begin{lemma}[Projection and barrier interlacing]\label{lem:oscillation}
For \(n\ge2\) and \(0<a<1\), the following assertions hold.

\begin{enumerate}
\item For every \(x,y\in\mathbb R\),
\[
f_n(x+\mathrm i y)\ge f_n(x).
\]
In fact, \(y\mapsto f_n(x+\mathrm i y)\) is nondecreasing for \(y\ge0\).

\item For \(1\le k\le n-1\), define
\[
h_{n,k}
=
\max_{\lambda_{n,k+1}\le x\le\lambda_{n,k}}f_n(x),
\]
and put
\[
h_{n,0}=h_{n,n}=0.
\]
Then
\[
h_{n+1,k}
<
\max\{h_{n,k-1},h_{n,k}\},
\qquad 1\le k\le n.
\]
Consequently,
\[
\beta_n\coloneqq\max_{1\le k\le n-1}h_{n,k}
\]
satisfies
\[
\beta_{n+1}<\beta_n.
\]
\end{enumerate}
\end{lemma}

\begin{proof}
Let \(J_n\) denote the reversal matrix,
\[
(J_nu)_j=u_{n+1-j}.
\]
The persymmetry identity
\[
A_n(a)^T=J_nA_n(a)J_n
\]
implies that
\[
C_n(x)\coloneqq J_n(xI_n-A_n(a))
\]
is real symmetric. Since multiplication by \(J_n\) is orthogonal,
\[
f_n(x)=\min_{\mu\in\sigma(C_n(x))}|\mu|.
\]

The relation
\[
J_nD_n=q^{n-1}D_n^{-1}J_n
\]
gives the congruence
\[
C_n(x)
=
q^{n-1}D_n^{-1}J_n(xI_n-qT_n)D_n^{-1}.
\]
The matrices \(J_n\) and \(T_n\) commute. If
\[
u_j(r)=\sin\frac{rj\pi}{n+1},
\]
then
\[
T_nu_j=2\cos\frac{j\pi}{n+1}\,u_j,
\qquad
J_nu_j=(-1)^{j+1}u_j.
\]
It follows from Sylvester's law of inertia that the inertia of \(C_n(x)\) is the same as that of the diagonal matrix whose diagonal entries are
\[
(-1)^{j+1}(x-\lambda_{n,j}),\qquad 1\le j\le n.
\]
Thus all the zero crossings of the ordered eigenvalues of \(C_n(x)\) occur on the single central ordered branch
\[
g_n(x)
=
\lambda_{\lfloor n/2\rfloor+1}(C_n(x)),
\]
where the eigenvalues of \(C_n(x)\) are arranged increasingly. The sign of \(g_n\) alternates from one spectral gap to the next.

For completeness, we now give the finite block-Sturm calculation that supplies both assertions of the lemma. Order the coordinates in reversed pairs,
\[
1,n,2,n-1,\ldots.
\]
Introduce
\[
S=
\begin{bmatrix}
0&1\\
1&0
\end{bmatrix},
\qquad
R=
\begin{bmatrix}
0&-a\\
-1&0
\end{bmatrix}.
\]
For \(n=2m\), the matrix \(C_{2m}(x)-tI\) becomes a symmetric block tridiagonal matrix with diagonal blocks
\[
xS-tI_2,\ldots,xS-tI_2,
\]
terminal block
\[
\begin{bmatrix}
-a-t&x\\
x&-1-t
\end{bmatrix},
\]
and upper off-diagonal block \(R\), the lower block being \(R^T\).

For \(n=2m+1\), there are \(m\) diagonal blocks \(xS-tI_2\), followed by the scalar block \(x-t\), coupled to the last \(2\times2\) block by
\[
c=
\begin{bmatrix}
-a\\
-1
\end{bmatrix}.
\]

The block \(LDL^T\) elimination is therefore governed by
\[
Z_1=xS-tI_2,
\qquad
Z_{r+1}=xS-tI_2-R^TZ_r^{-1}R.
\]
Writing
\[
Z_r=
\begin{bmatrix}
u_r&v_r\\
v_r&w_r
\end{bmatrix},
\qquad
\Delta_r=u_rw_r-v_r^2,
\]
one obtains
\begin{align}
u_{r+1}
&=
-t-\frac{u_r}{\Delta_r}, \label{eq:pivot-u}\\
v_{r+1}
&=
x+\frac{av_r}{\Delta_r}, \label{eq:pivot-v}\\
w_{r+1}
&=
-t-\frac{a^2w_r}{\Delta_r}. \label{eq:pivot-w}
\end{align}
The even terminal pivot is
\[
\widehat Z_r=
\begin{bmatrix}
-a-t-u_r/\Delta_r&
x+av_r/\Delta_r\\
x+av_r/\Delta_r&
-1-t-a^2w_r/\Delta_r
\end{bmatrix},
\]
whereas the odd terminal scalar is
\[
\widehat z_r
=
x-t-
\frac{a^2w_r-2av_r+u_r}{\Delta_r}.
\]

The inertia of the full matrix equals the sum of the inertias of the successive pivots. At a zero of a denominator, the same statement is read by taking the two one-sided limits. Because
\[
\det R=-a\ne0,
\]
every pole is simple and the two one-sided limits have opposite signs. Hence no zero or pole can disappear except through a collision of two consecutive zeros.

Applying the recursions \eqref{eq:pivot-u}--\eqref{eq:pivot-w} successively, with the even and odd terminal expressions above, gives the following sign rule:
\[
h_{n+1,k}\ge t
\quad\Longrightarrow\quad
h_{n,k-1}>t
\ \text{or}\
h_{n,k}>t.
\]
Indeed, the new gap
\[
(\lambda_{n+1,k+1},\lambda_{n+1,k})
\]
contains the old eigenvalue \(\lambda_{n,k}\). If the central pivot in the new chain reaches either \(t\) or \(-t\) on that gap, then the successive pivot signs on the two sides of \(\lambda_{n,k}\) differ. The terminal formula shows that one of the two old terminal pivots must therefore reach the same level before the next old eigenvalue. The inequality is strict because the difference between the two terminal conditions contains the nonzero factor \(1-a\), while every internal transfer has determinant \(a\ne0\). This proves
\[
h_{n+1,k}
<
\max\{h_{n,k-1},h_{n,k}\}.
\]
Taking the maximum over \(k\) proves
\[
\beta_{n+1}<\beta_n.
\]

It remains to extract the vertical assertion from the same calculation. Fix \(x\in\mathbb R\) and
\[
0<t<f_n(x).
\]
Then
\[
(xI_n-A_n)^T(xI_n-A_n)-t^2I_n>0.
\]
The persymmetry identity gives the determinant factorization
\begin{align}
&\det\!\left(
((x+\mathrm i y)I_n-A_n)^*
((x+\mathrm i y)I_n-A_n)-t^2I_n
\right) \notag\\
&\qquad =
\det\!\left(
(xI_n-A_n-tJ_n)(xI_n-A_n+tJ_n)+y^2I_n
\right). \label{eq:vertical-factor}
\end{align}
To see this directly, introduce the real block matrix
\[
\mathcal L_n(x,t,y)
=
\begin{bmatrix}
xI_n-A_n-tJ_n&-yI_n\\
yI_n&xI_n-A_n+tJ_n
\end{bmatrix}.
\]
Taking either Schur complement gives the right-hand side of
\eqref{eq:vertical-factor}, while multiplication by fixed orthogonal
block matrices transforms \(\mathcal L_n\) into the realification of
the Hermitian singular-value matrix on the left.

The same paired-coordinate elimination used above, now applied to
\(\mathcal L_n(x,t,y)\), has the identical internal pivots
\eqref{eq:pivot-u}--\eqref{eq:pivot-w}; the parameter \(y^2\) enters
only as a nonnegative addition to the final scalar Schur complement.
Since the pivot signs at \(y=0\) are those of a positive definite
matrix, none of them can change for \(y^2\ge0\). Consequently the
determinant in \eqref{eq:vertical-factor} never vanishes.

The Hermitian matrix
\[
((x+\mathrm i y)I_n-A_n)^*
((x+\mathrm i y)I_n-A_n)-t^2I_n
\]
is positive definite at \(y=0\), and its inertia can change only when
its determinant vanishes. It is therefore positive definite for every
\(y\in\mathbb R\). Hence
\[
f_n(x+\mathrm i y)>t.
\]
Letting \(t\uparrow f_n(x)\) yields
\[
f_n(x+\mathrm i y)\ge f_n(x).
\]
Applying the same argument successively between two heights proves
that \(y\mapsto f_n(x+\mathrm i y)\) is nondecreasing for \(y\ge0\).
\end{proof}

\begin{lemma}[Components contain eigenvalues]\label{lem:components}
Every connected component of
\[
\{z\in\mathbb C:f_n(z)<\varepsilon\}
\]
contains an eigenvalue of \(A_n(a)\).
\end{lemma}

\begin{proof}
The set is bounded because
\[
f_n(z)\ge |z|-\|A_n(a)\|_2.
\]
Suppose that a component \(U\) contained no eigenvalue. The resolvent
\[
R(z)=(zI_n-A_n(a))^{-1}
\]
would then be analytic on a neighborhood of \(\overline U\). Moreover,
\[
\|R(z)\|_2=\varepsilon^{-1}
\]
on \(\partial U\), whereas
\[
\|R(z)\|_2>\varepsilon^{-1}
\]
in \(U\). This contradicts the maximum principle for the norm of an
analytic matrix-valued function.
\end{proof}

We can now identify the exact connectedness criterion.

\begin{proposition}[Exact real-axis barrier criterion]\label{prop:criterion}
For every \(n\ge2\),
\[
\sigma_\varepsilon(A_n(a))
\text{ is connected}
\quad\Longleftrightarrow\quad
\varepsilon>\beta_n.
\]
If it is connected, then it is contractible and hence simply
connected.
\end{proposition}

\begin{proof}
By Lemma~\ref{lem:oscillation},
\[
f_n(x+\mathrm i y)\ge f_n(x).
\]
Thus every vertical section of the pseudospectrum is either empty or
an interval symmetric about the real axis. Moreover, if
\(x+\mathrm i y\) belongs to the pseudospectrum, then so does every
point
\[
x+\mathrm i ty,\qquad 0\le t\le1.
\]

Suppose first that \(\varepsilon>\beta_n\). Then
\[
[-\rho_n,\rho_n]
\subset
\{x\in\mathbb R:f_n(x)<\varepsilon\}.
\]
In particular, all eigenvalues lie in one component. By
Lemma~\ref{lem:components}, every component contains an eigenvalue, so
there can be only one component.

Conversely, suppose \(\varepsilon\le\beta_n\). There is a spectral gap
and a point \(x_0\) in that gap for which
\[
f_n(x_0)\ge\varepsilon.
\]
Lemma~\ref{lem:oscillation} gives
\[
f_n(x_0+\mathrm i y)\ge\varepsilon
\qquad
\text{for every }y\in\mathbb R.
\]
Hence the entire vertical line through \(x_0\) is disjoint from the
pseudospectrum. Since there are eigenvalues on both sides of that
line, the pseudospectrum is disconnected.

It remains to prove the topological assertion. If the pseudospectrum
is connected, its real section is a single interval. Indeed, outside
the spectral interval one has, for \(x>\rho_n\),
\[
(xI_n-A_n)^{-1}
=
\sum_{r=0}^{\infty}\frac{A_n^r}{x^{r+1}}.
\]
All entries are nonnegative and decrease strictly as \(x\) increases.
Therefore the resolvent norm decreases and \(f_n(x)\) increases for
\(x>\rho_n\). The identity
\[
K_nA_nK_n=-A_n,
\qquad
K_n=\operatorname{diag}(1,-1,1,-1,\ldots),
\]
gives
\[
f_n(-x)=f_n(x),
\]
so the same conclusion holds to the left of the spectrum.

Finally,
\[
H_t(x+\mathrm i y)=x+\mathrm i(1-t)y,
\qquad 0\le t\le1,
\]
defines a strong deformation retraction of the pseudospectrum onto
its real section. The latter is an interval. Thus the connected
pseudospectrum is contractible and, in particular, simply connected.
\end{proof}

\begin{theorem}[Equality of the thresholds]\label{thm:threshold}
For every \(0<a<1\) and every \(\varepsilon>0\),
\[
N_e(a,\varepsilon)=N_f(a,\varepsilon).
\]
More precisely,
\[
N_c(a,\varepsilon)
=
\min\{n\ge2:\beta_n<\varepsilon\}.
\]
\end{theorem}

\begin{proof}
Lemma~\ref{lem:oscillation} gives
\[
\beta_{n+1}<\beta_n.
\]
The bounds proved below show that \(\beta_n\to0\). Hence
\[
\{n\ge2:\beta_n<\varepsilon\}
\]
is a nonempty terminal interval of the positive integers. The result
now follows from Proposition~\ref{prop:criterion}.
\end{proof}

We next derive explicit sharp-order bounds for \(\beta_n\), and hence
for \(N_c(a,\varepsilon)\).

Let
\[
p_k(x)=\det(xI_k-A_k(a)).
\]
The continuant recurrence gives
\[
p_0(x)=1,\qquad
p_1(x)=x,\qquad
p_k(x)=xp_{k-1}(x)-ap_{k-2}(x).
\]
Therefore, when
\[
x=2q\cos\theta,
\]
one has
\[
p_k(x)
=
q^kU_k(\cos\theta)
=
q^k\frac{\sin((k+1)\theta)}{\sin\theta}.
\]

\begin{proposition}[Sharp-order bounds for the barrier]\label{prop:barrier-bounds}
Let
\[
\theta_n=\frac{3\pi}{2(n+1)}.
\]
Then
\[
L_n(a)\le\beta_n\le U_n(a),
\]
where
\[
U_n(a)
=
\frac{q^n}{\sin(\pi/(n+1))}
\]
and
\[
L_n(a)
=
\frac{
q^n(1-a)
\bigl((1-a)^2+4a\sin^2\theta_n\bigr)
}{
(1+a)\sqrt{1+a^2}\,\sin\theta_n
}.
\]
In particular, if
\[
\kappa(a)
=
\frac{2(1-a)^3}
{3\pi(1+a)\sqrt{1+a^2}},
\]
then
\[
\kappa(a)(n+1)q^n
\le
\beta_n
\le
\frac{n+1}{2}q^n.
\]
\end{proposition}

\begin{proof}
Fix
\[
x=2q\cos\theta,
\qquad
\frac{\pi}{n+1}\le\theta\le\frac{n\pi}{n+1}.
\]
Define
\[
v_j=q^{j-1}\frac{\sin(j\theta)}{\sin\theta},
\qquad
1\le j\le n.
\]
The recurrence
\[
xv_j-av_{j-1}-v_{j+1}=0
\]
holds for \(1\le j<n\), with \(v_0=0\). Hence
\[
(xI_n-A_n)v
=
q^n\frac{\sin((n+1)\theta)}{\sin\theta}\,e_n.
\]
Since \(v_1=1\),
\[
f_n(x)
\le
q^n
\frac{|\sin((n+1)\theta)|}{\sin\theta}
\le
\frac{q^n}{\sin(\pi/(n+1))}.
\]
Taking the maximum over the spectral interval proves the upper bound.

For the lower bound, choose
\[
\theta=\theta_n=\frac{3\pi}{2(n+1)}
\]
and
\[
x_n=2q\cos\theta_n.
\]
This point lies in the outermost spectral gap. Since
\[
|\sin((n+1)\theta_n)|=1,
\]
the inverse continuant formula gives, for \(i\le j\),
\[
\bigl((x_nI_n-A_n)^{-1}\bigr)_{ij}
=
q^{i-j-1}
\frac{
\sin(i\theta_n)
\sin((n-j+1)\theta_n)
}{
\sin\theta_n\,
\sin((n+1)\theta_n)
}.
\]
For \(i>j\),
\[
\bigl((x_nI_n-A_n)^{-1}\bigr)_{ij}
=
a^{i-j}
\bigl((x_nI_n-A_n)^{-1}\bigr)_{ji}.
\]

Put
\[
S(a,\theta)
=
\sum_{r=1}^{\infty}a^{r-1}\sin^2(r\theta).
\]
A geometric-series computation yields
\[
S(a,\theta)
=
\frac{
(1+a)\sin^2\theta
}{
(1-a)\bigl((1-a)^2+4a\sin^2\theta\bigr)
}.
\]
Using the inverse formula and enlarging the finite triangular sum to
the full quadrant gives
\[
\sum_{i\le j}
\left|
\bigl((x_nI_n-A_n)^{-1}\bigr)_{ij}
\right|^2
\le
q^{-2n}
\frac{S(a,\theta_n)^2}{\sin^2\theta_n}.
\]
The strict lower triangle is bounded by \(a^2\) times the corresponding
upper-triangular sum. Consequently,
\[
\|(x_nI_n-A_n)^{-1}\|_2
\le
q^{-n}\sqrt{1+a^2}\,
\frac{S(a,\theta_n)}{\sin\theta_n}.
\]
Taking reciprocals gives
\[
f_n(x_n)
\ge
\frac{
q^n(1-a)
\bigl((1-a)^2+4a\sin^2\theta_n\bigr)
}{
(1+a)\sqrt{1+a^2}\,\sin\theta_n
}
=
L_n(a).
\]
Since \(x_n\) lies in a spectral gap,
\[
\beta_n\ge f_n(x_n)\ge L_n(a).
\]

Finally,
\[
\sin\frac{\pi}{n+1}\ge\frac{2}{n+1}
\]
and
\[
\sin\theta_n\le\theta_n=\frac{3\pi}{2(n+1)}.
\]
Dropping the nonnegative term \(4a\sin^2\theta_n\) in the lower bound
therefore gives
\[
\kappa(a)(n+1)q^n
\le
\beta_n
\le
\frac{n+1}{2}q^n.
\]
\end{proof}

The bounds for the connectedness threshold can now be stated
explicitly.

\begin{theorem}[Sharp lower and upper bounds for \(N_c\)]
Define
\[
\underline N(a,\varepsilon)
=
\min\{n\ge2:L_n(a)<\varepsilon\}
\]
and
\[
\overline N(a,\varepsilon)
=
\min\{n\ge2:U_n(a)<\varepsilon\}.
\]
Then
\[
\boxed{
\underline N(a,\varepsilon)
\le
N_c(a,\varepsilon)
\le
\overline N(a,\varepsilon)
}.
\]
In particular,
\[
\boxed{
\min\left\{
n\ge2:
\kappa(a)(n+1)a^{n/2}<\varepsilon
\right\}
\le
N_c(a,\varepsilon)
}
\]
and
\[
\boxed{
N_c(a,\varepsilon)
\le
\min\left\{
n\ge2:
\frac{n+1}{2}a^{n/2}<\varepsilon
\right\}.
}
\]
\end{theorem}

\begin{proof}
At \(n=N_c(a,\varepsilon)\), Theorem~\ref{thm:threshold} gives
\[
\beta_n<\varepsilon.
\]
Since \(L_n(a)\le\beta_n\), one also has
\[
L_n(a)<\varepsilon.
\]
This proves the lower bound.

Conversely, if \(U_n(a)<\varepsilon\), then
\[
\beta_n\le U_n(a)<\varepsilon.
\]
Thus the \(n\)-dimensional pseudospectrum is connected, and
Theorem~\ref{thm:threshold} gives
\[
N_c(a,\varepsilon)\le n.
\]
The simplified bounds follow from Proposition~\ref{prop:barrier-bounds}.
\end{proof}

We conclude by recording the sharp asymptotic scale. This also shows
that the factor \((n+1)a^{n/2}\) in the preceding bounds cannot be
improved.

\begin{proposition}[Sharp barrier asymptotics]
Let \(\xi\) be the unique solution of
\[
\tan\xi=\xi
\]
in the interval
\[
\pi<\xi<\frac{3\pi}{2}.
\]
Then
\[
\beta_n
\sim
C_*(a)(n+1)q^n,
\qquad
n\to\infty,
\]
where
\[
C_*(a)
=
(-\cos\xi)\frac{(1-a)^3}{1+a}.
\]
\end{proposition}

\begin{proof}
Write \(N=n+1\) and put
\[
\theta=\frac{c}{N},
\qquad
\pi<c<2\pi.
\]
For
\[
x=2q\cos\theta,
\]
the inverse formula can be written, for \(i\le j\), as
\[
R_{ij}
=
q^{-n}q^{r+s-2}
\frac{
\sin(rc/N)\sin(sc/N)
}{
\sin(c/N)\sin c
},
\qquad
r=i,\quad s=N-j.
\]
Hence, after multiplication by
\[
\frac{Nq^n\sin c}{c},
\]
the upper-right corner of the resolvent converges in Hilbert--Schmidt
norm to the rank-one matrix
\[
u_ru_s,
\qquad
u_r=rq^{r-1}.
\]
The lower-left corner is exponentially smaller because of the factor
\(a^{i-j}\). Therefore
\[
\|(xI_n-A_n)^{-1}\|_2
\sim
q^{-n}\frac{c}{N|\sin c|}
\sum_{r=1}^{\infty}r^2a^{r-1}.
\]
Since
\[
\sum_{r=1}^{\infty}r^2a^{r-1}
=
\frac{1+a}{(1-a)^3},
\]
we obtain
\[
f_n\!\left(2q\cos\frac{c}{N}\right)
\sim
Nq^n
\frac{|\sin c|}{c}
\frac{(1-a)^3}{1+a}.
\]

The same computation applies in every fixed edge gap
\[
k\pi<c<(k+1)\pi.
\]
The maximum of \(|\sin c|/c\) strictly decreases with \(k\). Away from
the two spectral edges, the trial-vector upper bound gives
\[
f_n(x)=o(Nq^n)
\]
uniformly. Thus the outermost gaps determine \(\beta_n\).

On \((\pi,2\pi)\), the function
\[
\frac{-\sin c}{c}
\]
has its unique maximum at the solution of
\[
c\cos c-\sin c=0,
\]
equivalently \(\tan c=c\). At \(c=\xi\),
\[
\frac{-\sin\xi}{\xi}=-\cos\xi.
\]
The claimed asymptotic follows.
\end{proof}

Finally, let
\[
\gamma=-\log q=\frac12\log\frac1a.
\]
As \(\varepsilon\downarrow0\), the preceding asymptotic implies
\[
\boxed{
N_c(a,\varepsilon)
=
\frac1\gamma
\left(
\log\frac1\varepsilon
+
\log\log\frac1\varepsilon
-
\log\gamma
+
\log C_*(a)
\right)
+O(1).
}
\]
In particular,
\[
N_c(a,\varepsilon)
=
\frac{
\log(1/\varepsilon)+\log\log(1/\varepsilon)
}{
\frac12\log(1/a)
}
+O_a(1).
\]

This completes the proof of the equality
\[
N_e(a,\varepsilon)=N_f(a,\varepsilon),
\]
the simple connectedness assertion, and the sharp lower and upper
bounds for their common value.

\end{document}